%! The Normal Distribution and z-Scores — Unit 1, Lesson 1.7 — Warm-Up (Answer Key)
% ONE PAGE, blank and key. Three items.
% Mirrors the blank exactly apart from the package swap, the pageheader, and the answers.
% NO teachernote here — teacher prose lives in the lesson plan.
\documentclass[10pt]{article}
\usepackage{apstats-article}
\usepackage{apstats-key}   % pulls in -boxes; do NOT also load -boxes

\begin{document}

\pageheader{Unit 1, Lesson 1.7}{Warm-Up --- Answer Key}

\vspace{0.15in}

\begin{enumerate}[leftmargin=1.7em, itemsep=0.19in]

  \item \textbf{From last lesson.} Here are five-number summaries for the daily commute times,
        in minutes, of two groups of workers.

        \vspace{4pt}
        \renewcommand{\arraystretch}{1.4}
        \begin{tabularx}{\linewidth}{@{} l Y Y Y Y Y @{}}
        \rowcolor{sky}
        \textbf{Group} & \textbf{Min} & \textbf{Q1} & \textbf{Med} & \textbf{Q3} & \textbf{Max} \\
        \hline
        \textbf{Downtown} & 8 & 14 & 19 & 26 & 41 \\
        \textbf{Suburban} & 15 & 22 & 25 & 30 & 38 \\
        \end{tabularx}
        \vspace{5pt}

        IQR for Downtown: \ans{12 min} \qquad IQR for Suburban: \ans{8 min}

        \vspace{4pt}
        Whose \emph{middle half} of commute times is more spread out? \ans{Downtown (12 vs.\ 8)}

  \item A set of quiz scores has a mean of \textbf{20} points and a standard deviation of
        \textbf{3} points.

        \vspace{5pt}
        \hspace{1.0em}\textbf{(a)} How many points \emph{above} the mean is a score of 26?
        \ans{6 points}
        \vspace{5pt}

        \hspace{1.0em}\textbf{(b)} That many points is how many 3-point steps? \ans{2 steps}
        \vspace{5pt}

        \hspace{1.0em}\textbf{(c)} A score of 17 is how many 3-point steps from the mean, and on
        which side? \ans{1 step, below the mean}

  \item In a set of \textbf{400} measurements, \textbf{360} of them fall inside a certain range
        of values.

        \vspace{5pt}
        \hspace{1.0em}\textbf{(a)} What percent of the measurements is that? \ans{90\%}
        \vspace{5pt}

        \hspace{1.0em}\textbf{(b)} What percent fall \emph{outside} that range? \ans{10\%}
        \vspace{5pt}

        \hspace{1.0em}\textbf{(c)} The measurements pile up evenly on both sides of the middle,
        so the leftovers split evenly between the low end and the high end. What percent fall
        past the \emph{high} end? \ans{5\%}

\end{enumerate}

\end{document}
